re-tee

It takes two arguments: the process id of a running program and a
file name. It adds code to that program that
copies to the named file all output that the program writes to its
standard output file descriptor. In this way it
works like ``tee,'' which passes output along to its own standard out
while also saving it in a file. The motivation
for the example program is that you run a program, and it starts to
print copious lines of output to your screen, and
you wish to save that output in a file without having to re-run the program.

\lstinputlisting{examples/retee.C}
